\documentclass[12pt]{article}
\usepackage[utf8]{inputenc}
\usepackage[T1]{fontenc}
\usepackage{hyperref}
\title{Saba Language Resources for K-12 Curriculum Development}
\author{Compiled by Claude Sonnet 5.5 (claude-sonnet-5-5) for Andr\'as Kornai}
\date{}
\begin{document}
\maketitle

\section{Introduction}
Saba (ISO 639-3 saa, Glottocode saba1276) is an East Chadic (Sokoro group) language of south-central Chad, spoken in Sorki canton, Chinguil sub-prefecture, in the Melfi region. Ethnologue treats Jelkung as a synonym, while Blench (2006) treats it as a separate language. Estimates are about 1,300 to 1,500 speakers. Saba is not related to the Central Sudanic Sara and Kaba languages, which are easily confused with it. Only academic and survey material turned up. I found no primers, school books or community programmes.

\section{Primary Educational Resources}
None found, although the reference description says Saba is taught as a subject in some schools. I found no evidence of the materials used.

\section{Linguistic and Reference Materials}
\begin{itemize}
\item Jungraithmayr, \emph{Preliminary studies of Saba, an Eastern Chadic language in the Melfi region of Chad} (\emph{Afrika und \"Ubersee}). It covers verb morphology. \url{https://journals.sub.uni-hamburg.de/hup1/afrikaunduebersee/article/view/221}
\item Dakouli, Maa\ss{} and Toomey (1996), a 35-page SIL sociolinguistic survey of Saba, catalogued by Glottolog. \url{https://glottolog.org/resource/reference/id/157642}
\item ASJP Saba wordlist (compiled by Julia Bischoffberger from the 1996 survey). \url{https://asjp.clld.org/languages/SABA}
\item Wikipedia article on Saba. \url{https://en.wikipedia.org/wiki/Saba_language}
\item Wikipedia article on Jelkung. \url{https://en.wikipedia.org/wiki/Jelkung_language}
\item Endangered Languages Project entry. \url{https://endangeredlanguages.com/lang/2592}
\end{itemize}

\section{Religious and Community Resources}
None found.

\section{Digital and Audio Resources}
None found beyond the ASJP wordlist above.

\section{Curriculum Development Context}
Saba has a very small speaker population and no known standard orthography or pedagogical materials. Curriculum work would have to start from the survey and the Jungraithmayr grammar notes. It would need new documentation with the community, an agreed alphabet, and coordination with Chadian education authorities. The reported school teaching could be a starting point for finding teachers, but this needs confirming locally.
\end{document}